% ===== §6 The Fourth Cause: The Order-Problem of Opaque Subjects =====
\section{The Fourth Cause: Opaque Subjects and the Order-Problem}
\label{p17:sec:ir}

\noindent
The three internal causes were traced within the dyad, as though the two subjects stood alone. They do not. Around them runs a field of powers, and through them runs the question those powers pose: how can two beings who cannot enter one another, who answer to no authority above them, and who act always under uncertainty, nonetheless build between them an order---a trust, a coordination, a peace? This is the question of \emph{international relations}; and the burden of this section is that, properly reconstructed, it is also, and more purely, the question of the intimate bond. The reconstruction is delicate, because the discipline of IR carries within it assumptions that the generative-relational account must reject, and others it urgently needs. The section separates them.

\subsection{Reconstructing the object of IR}
\label{p17:sub:ir-redefine}

What is the theory of international relations a theory \emph{of}? The realist tradition gives one answer: it is the theory of how states, as like units in a condition of anarchy, pursue survival under the imperative of \emph{relative power}, so that the fundamental currency of the system is the relative gain---what one unit has only insofar as another lacks it.\footnote{For the structural-realist core---anarchy, like units, survival, and the primacy of relative over absolute gains---see \textcite{waltz1979}; for its offensive variant, \textcite{mearsheimer2001}.} On this answer, the master concepts are the balance of power, the security dilemma in its material form, and the alliance as an instrument of relative advantage. And on this answer, IR is useless to us---worse than useless, for its central assumption is precisely the one the generative-relational account identifies as the mark of a \emph{bad} cycle. A bond in which value is reckoned in relative gains, in which one subject has worth only insofar as the other lacks it, is not a hard case of intimacy but a relation that has already decayed into extraction (zero or negative holonomy). To import the realist currency into the dyad is not to analyse intimacy but to assume its corruption.

But realism is not IR. Stripped of its particular and contestable hard core, the discipline has a deeper and more general object, which its own constructivist and English-School traditions have always pursued: the problem of how \emph{opaque actors build order under uncertainty without a central authority}. Three conditions define this object, and not one of them mentions relative power. The actors are \emph{opaque}---they cannot read one another's interiors and must infer them. They act under \emph{uncertainty}---no actor knows another's true disposition or future conduct. And they stand in \emph{anarchy} in the strict and non-pejorative sense: there is no authority above them to enforce agreements or adjudicate disputes.\footnote{For anarchy as a structure whose meaning is constituted by the practices of its actors rather than given---``anarchy is what states make of it''---see \textcite{wendt1992}; for the conception of an \emph{international society} in which order arises from shared norms and institutions rather than from power alone, \textcite{bull1977}.}

\begin{claim}[The reconstructed object of IR]
The deep object of international relations is not the pursuit of relative power among states but the more general problem of \emph{how opaque subjects, unable to enter one another directly, build trust and order under uncertainty and in the absence of any authority above them}. So defined, the intimate dyad is itself an instance of the IR problem---and a purer instance than the international one, because the opacity of subjects is here \emph{ontological} rather than merely practical: it is not that the other's interior is hard to observe, but that it is constitutively unsymbolisable, circled by a vacancy no observation could fill.
\end{claim}

\noindent
This is the pivot on which the section turns. The intimate bond satisfies all three defining conditions more exactly than any state system: no sword can be borrowed to enforce a vow between two persons (the anarchy is complete), no future fidelity can be known in advance (the uncertainty is total), and no lover can read the beloved's heart directly (the opacity is ontological, not technical).\footnote{On the impossibility of borrowing an external sword to bind an intimate vow---the constitutive anarchy of the dyad---see the account of the non-binding vow and relational self-legislation in \textcite{paperxiii}.} What IR studies under the cover of states, the intimate bond exhibits in its pure form. And the most developed contemporary statement of this reconstructed IR comes, fittingly, from a relational theory of world politics that the next subsections will both draw on and demarcate.

\subsection{Power is ineliminable within the dyad}
\label{p17:sub:ir-internal}

The first temptation, once the realist currency is rejected, is to over-purify: to imagine the dyad's interior as a power-free zone, a clearing of pure generativity into which power intrudes only from outside. This temptation must be refused, and refusing it is the precondition of everything that follows. \emph{Power is internal to the dyad, ineliminably.}

The reason is ontological and was given in the introduction. The subjects are symbolic subjects---subjects constituted in and through the symbolic order---and wherever there is a symbolic relation there is already an asymmetry: whose word counts as the neutral description and whose as the special pleading; whose grammar of trust (\xref{p17:sec:encoding}) is treated as the default and whose as the deviation; whose desire sets the rhythm the other must read. These are not intrusions of external power; they are the micro-physics of any relation between symbolic beings. There is no coupling of two subjects so loving that it contains no asymmetry, because to be a symbolic subject in relation is already to occupy a position, and positions are not equal.

It follows that the health of a bond cannot consist in the \emph{absence} of internal power. It must consist in something else:

\begin{claim}[Health is not the absence of power]
A healthy intimate bond is not one from which internal power has been eliminated---that is impossible for symbolic subjects---but one in which power is continuously \emph{transformed by the generative logic} rather than allowed to solidify into standing domination. Internal power is not a pathology to be cured but a permanent condition to be metabolised. The pathology is not power; it is its freezing.
\end{claim}

\noindent
This reframing does real work later. It means the internal causes of \xref{p17:sec:encoding}--\xref{p17:sec:phase} were never power-free either: the default grammar, the clear or noisy channel, the curvature of the manifold are all already shot through with micro-asymmetries. And it means that the relational theory of world politics, to which we now turn for resources, must be taken up critically---for it, too, risks an over-sweet picture of relation as pure mutual empowerment.

\subsection{The continuous gradient, and regime bleeding}
\label{p17:sub:ir-gradient}

If power is internal as well as external, then the picture of two discrete zones---a generative inside and a power-laden outside---is too crude. What the dyad actually inhabits is a \emph{gradient}: a continuous field in which the concentration of power-logic rises as one moves outward from the dyadic core, through near kin, to extended family and clan, to the institutions and scripts of the wider society. At the core, power is present but dilute, held in solution by the generative logic; further out, it grows denser, less metabolised, more apt to crystallise into standing structure. There is no sharp boundary between a power-free interior and a power-laden exterior, only a rising concentration along a continuum.

This gradient picture is what makes intelligible the most insidious of the external causes, and the one the paper marks as a central original contribution: \emph{regime bleeding}. Because the inside and the outside are continuous rather than walled off, the power-logic dense at the periphery can \emph{seep inward} along the gradient and contaminate the dyadic core---can install, at the heart of a coupling that should run on the non-zero-sum logic of relational wealth, the relative-gains calculus proper to the periphery.

\begin{claim}[Regime bleeding]
Because power varies continuously from the dyadic core outward rather than dividing into walled zones, the dense power-logic of the periphery can seep inward and contaminate the core, converting a non-zero-sum generative coupling into a relative-gains calculus. This is the most insidious of the external causes precisely because no boundary is breached---the contamination travels a continuous gradient---and because it can corrupt a core that is, in itself, entirely healthy. ``To manage my family, you must prove you are worth it'': the sentence imports, into the heart of the bond, an accounting that the bond cannot survive. A healthy dyad can be hollowed from the periphery inward without any internal fault whatever.
\end{claim}

\noindent
Regime bleeding (cause 15 of the etiology) is the hinge between the internal and external causes, and it explains why the dyad cannot simply be sealed against the field. The seepage is possible only because the boundary is semi-permeable---and the boundary \emph{must} be semi-permeable, because a dyad sealed absolutely against the field would be a dyad cut off from the external connections on which, as the praxis will argue (\xref{p17:sec:praxis}), its own internal balance partly depends. The same permeability that lets in the noise lets in the air. This is the first appearance of a structural cause (cause 17) that the cosmological sections will take up: the semi-permeable boundary is not a defect to be sealed but a permanent condition to be managed.

\subsection{Estrangement as the interaction-failure of opaque subjects}
\label{p17:sub:ir-failure}

With the reconstructed object in hand, the four lenses of the etiology can be seen to converge. \emph{Surechigai}, viewed through the reconstructed IR, is the characteristic interaction-failure of two opaque subjects building order without an arbiter: it is \emph{misperception}, \emph{mistranslation}, \emph{withdrawal}, and the failure of \emph{return}, named all at once. Each of the prior lenses described one facet of a single thing. The trust-grammar mismatch (\xref{p17:sec:encoding}) is mistranslation: the IR literature on \emph{misperception} describes exactly this---the systematic divergence between a signal's intended and received meaning, arising not from bad faith but from the differing cognitive schemata of sender and receiver, and unknown, often, to both.\footnote{For misperception as a systematic, schema-driven divergence between intended and received signals---a cognitive rather than a power-theoretic phenomenon---see \textcite{jervis1976}.} The unattributable error and predictive withdrawal (\xref{p17:sec:prediction}) is the \emph{security dilemma} in its cognitive, not its material, form: a defensive disinvestment, undertaken purely to stanch one's own unmanageable error, that the other cannot but read as hostility, and answers with disinvestment of their own.\footnote{The security dilemma admits two readings. Its realist version (\cite{mearsheimer2001}) turns on relative power and is rejected here; its cognitive version turns on uncertainty and the misreading of defensive behaviour as threat, and requires no concept of relative power at all. It is only the latter, cognitive version that the present argument employs.} And the arrest of phase (\xref{p17:sec:phase}) is the failure of \emph{return}---the loop that, in IR's terms, never reaches the cooperative equilibrium that repeated interaction was supposed to build, because each round's unattributable defection resets it.

This convergence is the analytic payoff of the reconstruction. It shows that the four causes are not a list but a cascade, and that the reconstructed IR is the framework in which the cascade can be seen whole: opaque subjects, predicting one another across an unbridgeable interior, mistranslate (cause-layer one), cannot attribute the resulting error (cause-layer two), withdraw to protect themselves, and so fail to close the loop that would have generated value (cause-layer three)---all of it under no authority that could arbitrate, and all of it liable to contamination from a periphery dense with power (cause-layer four). The reconstructed IR does not add a fourth cause beside three others; it is the stage on which all four play out.

\subsection{Isomorphism, and the rupture: the two-level teleology}
\label{p17:sub:ir-rupture}

It remains to state, with precision, both how far the intimate bond is \emph{like} a state system and exactly where the likeness breaks---for the break is the deepest structural finding of the section, and it governs everything in the chapters on means that follow.

The likeness is real. Both the dyad and the state system are populations of opaque actors building order under uncertainty without an arbiter; both face misperception, the cognitive security dilemma, the problem of credible commitment, the temptation to withdraw. The strategies that opaque actors have devised for surviving such a condition---the subject of the next section---therefore transfer, with care, from the one domain to the other. To this extent the isomorphism holds, and IR earns its place in the argument.

But the likeness breaks at a single, decisive point, and the break is not one of degree but of \emph{teleological order}. For a state, in the realist understanding, survival is the terminal end: the state survives in order to survive, and there is no further good, beneath survival, that survival serves and might betray. The intimate bond is not like this. Its survival is not terminal but \emph{instrumental}: the two contend with power, balance, defend, and endure \emph{in order to protect} something that is not itself a matter of power at all---the dynamics of love, the generativity of the real coupling, the non-power order that is the bond's actual end. The dyad is a \emph{two-level} system where the state system is single-level.

\begin{claim}[The two-level rupture]
The intimate bond is isomorphic to a state system in its surface problem---opaque actors building order under uncertainty without an arbiter---but ruptures from it in teleological order. For the state, survival is the terminal end. For the dyad, survival of the power-order is merely the \emph{means} to a further and non-power end: the generativity of the real coupling it exists to protect. The dyad is two-level; the state system, single-level. And the means has an intrinsic tendency to damage the end---a defence mounted to protect the love can import the very relative-gains logic the love cannot survive---so that the dyad must do something a state never must: win its security \emph{without} letting the winning corrupt the thing secured.
\end{claim}

\noindent
This is the precise content of the claim, ventured loosely in the introduction, that the intimate relation faces, on certain axes, a problem \emph{more complex} than the international one, even as it is incomparably smaller in scale. The greater complexity is not a matter of more variables. It is the two-level structure itself: a state optimises a single objective (survival), whereas a dyad must pursue an instrumental objective (security against power) that is in standing tension with the terminal objective it serves (the generativity of love), under the constant risk that the means will consume the end. A small system with two incommensurable levels can be harder to govern than a large one with a single level. This is why the relational theory of world politics, for all its affinity with the present account, cannot simply be adopted: it rightly sees relation as a source of power that is \emph{power-to} rather than \emph{power-over}---a co-empowering, mutually generative capacity rather than a capacity for domination---but it does not equip itself with the holonomic criterion that distinguishes a good relational cycle from a bad one, nor with the account of the real that explains why the bond's terminal end cannot be reached by the very relational means that secure it.\footnote{For the relational theory of world politics---relation as itself a form of power, power as more often \emph{power-to} (co-empowerment) than \emph{power-over} (domination), and the \emph{yin--yang} dynamic of mutual completion as a generative process irreducible to either party---see \textcite{qin2018}. The present account draws on this conception of power-to while marking three points of divergence: it adds the holonomic criterion separating good from bad relational cycles (\cite{paperix}); it grounds relation in the Lacanian real and its unsymbolisable remainder, where the relational theory rests on a Confucian grammar of reciprocity; and it locates, at \xref{p17:sec:balance}, an ontological account of why such relation must continually generate a counter-force against its own freezing.}

The rupture sets the agenda for the second of the paper's three questions. The dyad must defend the conditions of its love against the field of power; the tools for that defence are the reconstructed strategies of IR; but those tools operate at the level of means, and the love they defend lives at the level of an end they cannot touch. To use them well is therefore to know exactly what they can and cannot reach---which is the burden of the next section, on the political economy of conditions, and its successor, on the inventory of strategies and the teleological filter that decides which of them a bond may use without destroying what it defends.
